%=============================================================================!
% 16.S  WHAT WE NOW HAVE                                                      !
%=============================================================================!
\unnumbered{What we now have}

The radial distribution function counts the partners at each distance
against an even scatter, a histogram of pair distances out to half the
cell that tends far away to $1 - S(0)/N$ in a fixed cell, $1 - 1/N$ for
atoms placed at random; with the pair potential and the kinetic energy it
gives the energy and the pressure, and with its short memory it converges
several times sooner than the pressure. Its transform is the structure factor
that diffraction measures, computed from the positions on the cell's own
wave vectors where the cut at half the cell would spoil it, and tending
to $\rho\kB T\kappa_T$ at long wavelengths.
Bonds are assigned from
distances with a band of hysteresis and a least lifetime, stored as
whole-number keys whose sorted lists give the bonds formed and broken;
molecules are the connected components of their graph, species their
formula and a hash of their bonding; reactions are changes of the
molecules that atoms belong to, and rates are counts over exposures with
the error $\sqrt n$; the ratio of the rate constants of a reaction and
its reverse is an equilibrium constant that statistical mechanics
predicts. Observers do this as the run goes, at about a tenth of the cost
of a cheap force evaluation.
The mean squared displacement grows as $2dDt$, from unwrapped
paths with the centre of mass removed, and equals twice the integral of
the velocity autocorrelation function weighted by $t - s$, so $D$ is also
a Green-Kubo integral; the shear viscosity is another, and its value,
$(1.853 \pm 0.038)\times10^{-4}$\,\si{\pascal\second} for liquid argon,
lies 1.6 combined errors from the one that Chapter~15 inferred from the
change of $D$ with the size of the box.
Independent ions conduct as $\sum_iq_i^2D_i/V\kB T$,
and the Haven ratio measures how far correlated ions fall short of it;
in layers, $D$ becomes a tensor; and hops between sites give $D =
fk_{\mathrm h}a^2/4$, with a correlation factor $f$ that friction brings to 1.
The Fourier transform, computed in $N_{\mathrm s}\log_2N_{\mathrm s}$
operations for $N_{\mathrm s}$ samples, turns the velocity
autocorrelation function into the vibrational density of states, whose
peaks are the normal modes, whose width grows by $\gamma/2\pi$ under
friction, whose value at zero holds $D$, and which frames a stride
$\Delta t$ apart show only up to $1/2\Delta t$.
A histogram of any
coordinate gives its free energy, which counts arrangements as well as
energies, with error bars that must come from independent runs when the
barriers are crossed rarely.

Chapter~17 runs these analyses on the simulations they are meant for:
ASE and LAMMPS with a learned potential, lithium in graphite, and the
trajectory gallery of every fault met so far.
